\bigskip

\textbf{Question 1}

\medskip

$f(x)=\left (2+5x\right ) \e^{3x}$ donc :
\begin{align*}
f'(x) &= 5\times \e^{3x} + \left (2+5x\right )\times 3\e^{3x} \\
&= \left (5+6+15x\right ) \e^{3x} \\
&= \left (11+15x\right )\e^{3x}
\end{align*}
La bonne réponse est \textbf{C}.

\bigskip

\textbf{Question 2}

\medskip

La solution générale d'une équation différentielle d'ordre 1 est la somme d'une solution particulière et de la solution générale de l'équation différentielle sans second membre associée.

\begin{list}{\textbullet}{}
\item L'équation sans second membre associée est: $y'=-3y$.

D'après le cours, on sait que les solutions de cette équation sont les fonctions $y_0$ définies par $y_0(x)= k\e^{-3x}$ où $k$ est un réel quelconque.

\item On cherche une solution particulière de $(E)$ sous forme de constante $a$. \\
Donc $a'=-3a+5$. Or $a'=0$, donc $-3a+5=0$ donc $a=\dfrac{5}{3}$.
\end{list}

Les fonctions solutions de l'équation différentielle $(E)$ sont les fonctions $y$ définies sur $\R$ par $y(x)= k \e^{-3x} +\dfrac{5}{3}$ où $k$ est un réel quelconque.

\bigskip

\textbf{Question 3}

\medskip

On veut déterminer la forme exponentielle du nombre complexe 
$z = - 6\sqrt 3 + 6\text{i}.$

\begin{list}{\textbullet}{}
\item Le module de $z$ est $|z| = \ds\sqrt{\left(-6\sqrt{3}\right )^2 + 6^2 } = \ds\sqrt{108+36} = \ds\sqrt{144}=12$.

\item Un argument $\theta$ de $z$ vérifie $\cos\left (\theta\right )=\dfrac{-6\sqrt{3}}{|z|} = \dfrac{-6\sqrt{3}}{12} = -\dfrac{\sqrt{3}}{2}$ et $\sin\left (\theta\right )=\dfrac{6}{|z|} = \dfrac{6}{12} = \dfrac{1}{2}$.

$\theta = \dfrac{5\pi}{6}$ vérifie ces égalités.
\end{list}

La forme exponentielle du nombre complexe 
$z = - 6\sqrt 3 + 6\text{i}$
est donc $z = 12 \e^{\frac{5\pi}{6}\i}$.

\bigskip

\textbf{Question 4}

\medskip

\begin{list}{\textbullet}{Graphiquement, on voit que:}
\item sur l'intervalle $]-\infty\,,\, -2[$, la courbe $\mathcal{C}$ est au-dessus de la droite $\Delta$;
\item sur l'intervalle $]-2\,,\, 4[$, la courbe $\mathcal{C}$ est en dessous de la droite $\Delta$;
\item sur l'intervalle $]4\,,\, +\infty[$, la courbe $\mathcal{C}$ est au-dessus de la droite $\Delta$.
\end{list}

Sur l'intervalle $[-2\,,\,4]$, la courbe est en dessous de la droite.

Soit $\mathcal{A}_1$ l'aire comprise entre la droite $\Delta$, l'axe des abscisses et les droites verticales d'équations $x=-2$ et $x=4$.\\
Soit $\mathcal{A}_2$ l'aire comprise entre la courbe $\mathcal{C}$, l'axe des abscisses et les droites verticales d'équations $x=-2$ et $x=4$.

L'aire hachurée est égale à $\mathcal{A}_1-\mathcal{A}_2$.

\begin{list}{\textbullet}{}
\item $\mathcal{A}_1$ est l'aire d'un trapèze rectangle de petite base 3, de grande base 9, et de hauteur 6, donc\\
 $\mathcal{A}_1=\dfrac{\left (3+ 9\right )\times 6 }{2}=\dfrac{72}{2}=36$.
\item $\mathcal{A}_2 = \ds\int_{-2}^{4}f(x) \d x$.

La fonction $f$ a pour primitive la fonction $F$ définie par $F(x)=\dfrac{1}{2}\times \dfrac{x^3}{3}+x$ soit $F(x)=\dfrac{x^3}{6}+x$.

$\ds\int_{-2}^{4}f(x) \d x 
= F(4)-F(-2)
= \left ( \dfrac{4^3}{6} + 4 \right ) - \left ( \dfrac{(-2)^3}{6} + (-2)\right )
= \dfrac{64}{6}+ 4 + \dfrac{8}{6} + 2
= \dfrac{72}{6}+6 = 12+6=18$.
\end{list}

L'aire hachurée vaut donc, en unités d'aire: $36 - 18 = 18$.

\textbf{Correction alternative :}

On peut calculer l'aire de la partie colorée en faisant :
$\ds\int_{-2}^{4}(g(x)-f(x)) \d x.$

Comme :
$g(x)-f(x) = x + 5 - \left(\dfrac12 x^2 +1\right) = -\dfrac12 x^2 + x + 4$

Alors :
\begin{align*}
\ds\int_{-2}^{4}\left(-\dfrac12 x^2 + x + 4\right) \d x &= \left[-\dfrac{1}{6}x^3 + \dfrac{1}{2}x^2 + 4x\right]_{-2}^{4}\\
&= \left(-\dfrac{1}{6} \times 4^3 + \dfrac{1}{2} \times 4^2 + 4 \times 4\right) - \left(-\dfrac{1}{6} \times (-2)^3 + \dfrac{1}{2} \times (-2)^2 + 4 \times (-2)\right) \\
&= \dfrac{40}{3} + \dfrac{14}{3} \\
&= 18 \text{ u.a.}
\end{align*}

\bigskip
